\section{Анализ роботоспособности системы мультиобъектного отслеживания}

Далее будет проведена проверка роботоспособности полученной системы мультиобъектного отслеживания. Проверка представляет из себя анализ основных метрик, которые применяются в системах такого рода и сравнение этих метрик с аналогичными технологиями. Для их получения рассматривается видеоматериал с различными условиями: освещение, погода, плотность трафика, окружающая обстановка. 

\subsection{Проверка работоспособности системы МОТ в различных условиях}

Для проверки работоспособности разработанной системы мультиобъектного отслеживания необходимо провести ряд экспериментов. В качестве экспериментов будут выступать видео с различными условиями:

\begin{itemize}
    \item По освещенности --- рассматриваются видео снятые в светлое время суток, в темное, а также с переходами из темного в светлое (из тоннеля) и из светлого в темное (в тоннель);
    \item По погодным условиям --- рассматриваются видео с ясной, пасмурной, туманной и заснеженной погодой;
    \item По трафику --- рассматриваются видео с не плотным, средним и плотным трафиками;
    \item По окружению --- рассматриваются видео в гороской среде, на трассе и извилистой дороге (серпантине).
\end{itemize}

Тестирование и сравнение параметров всех возможных вариантов для каждой из систем МОТ является довольно трудоемкой и масштабной задачей, так как в результате получим более тысячи видео.

Источниками видеоматериалов являются наборы данных, предназначенные для автономного вождения, такие как MOTChallenge(MOT17, MOT20), BDD100K, KITTI~\cite{cao2022observation}. Метрики, по которым будет проводится сравнение перечислены в части~\ref{subsec:metrics}. 
Таблица с полученными данными приведена в приложении~\ref{tab::mot_video_results}.

Далее рассмотрим сводные таблицы и графики, полученные в результате анализа экспериментальных данных. 

В таблице~\ref{tab:average_results} представлены среднее значение метрик, полученных по итогам тестирования. В ней собраны метрики со всех проведенных тестов.

\begin{table}[h!]
    \centering
    \caption{Сводная таблица средних значений метрик для системы МОТ во всех тестах}
    \label{tab:average_results}
    \begin{tabular}{|l|c|c|c|c|c|c|}
        \hline
        \textbf{Алгоритм} & \textbf{MOTA} & \textbf{IDF1} & \textbf{IDS} & \textbf{HOTA} & \textbf{FAF} & \textbf{FPS} \\
        \hline
        ByteTrack & 0.803 & 0.773 & 37.6 & 0.631 & 0.18 & 30 \\ \hline 
        \textbf{DeepSORT+EKF} & 0.811 & 0.811 & 26.3 & 0.661 & 0.052 & 45 \\ \hline
        DeepSORT+LKF & 0.767 & 0.706 & 44.2 & 0.548 & 0.076 & 47 \\ \hline
        FairMOT & 0.695 & 0.694 & 42.2 & 0.545 & 0.076 & 22 \\ \hline
        OC-SORT & 0.799 & 0.779 & 34.5 & 0.580 & 0.068 & 48 \\ \hline
    \end{tabular}
\end{table}

Как видно из таблицы~\ref{tab:average_results}, полученная система МОТ (DeepSORT+EKF) полностью удовлетворяет требованиям, которые перечислены в пункте 2 в части~\ref{sec:task}.

Далее перейдем к тестам со интересующими нас условиями, такими как: плохое освещение, дождь, туман, плотный траффик и тд. Данные, которе были получены по результатам тестирования представлены в таблице~\ref{tab:specific_results}

\begin{table}[h!]
    \centering
    \caption{Сводная таблица средних значений метрик для системы МОТ в тестах с плохими условиями}
    \label{tab:specific_results}
    \begin{tabular}{|l|c|c|c|c|c|c|}
        \hline
        \textbf{Алгоритм} & \textbf{MOTA} & \textbf{IDF1} & \textbf{IDS} & \textbf{HOTA} & \textbf{FAF} & \textbf{FPS} \\
        \hline
        ByteTrack & 0.737 & 0.737 & 38.1 & 0.587 & 0.066 & 50 \\\hline
        \textbf{DeepSORT+EKF} & 0.809 & 0.79 & 36.4 & 0.659 & 0.12 & 35 \\\hline
        DeepSORT+LKF & 0.695 & 0.694 & 44.7 & 0.545 & 0.077 & 47 \\\hline
        FairMOT & 0.693 & 0.692 & 42.7 & 0.543 & 0.077 & 22 \\\hline
        OC-SORT & 0.727 & 0.726 & 35.0 & 0.577 & 0.068 & 48 \\\hline
    \end{tabular}
\end{table}

Исходя из полученных результатов можно сделать вывод, что система показывает неплохие результаты в неблагоприятных условиях. Полученная система МОТ (DeepSORT+EKF) полностью удовлетворяет требованиям, которые перечислены в пункте 3 в части~\ref{sec:task}. 

Далее, на рисунках~\ref{img::mota_weather}-~\ref{img::mota_light} представлены графики сравнения параметров системы в различных условиях:

\img[!h]{img::mota_weather}{MOTA_weather.png}{MOTA при различных погодных условиях}{1}

\img[!h]{img::idf1_weather}{IDF1_weather.png}{IDF1 при различных погодных условиях}{1}

\img[!h]{img::mota_traffic}{MOTA_traffic.png}{MOTA при разном уровне траффика}{1}

\img[!h]{img::mota_light}{MOTA_light.png}{MOTA при различных погодных условиях}{1}

Также рассмотрим как зависит ошибка предсказания позиции от системы. Сравнительный график представлен на рисунке \ref{img::error}.

\img[!h]{img::error}{extended_position_error_comparison.png}{Ошибка предсказания позиции от системы}{1}

DeepSORT+EKF показывает лучший результат в обоих сценариях. Улучшение EKF относительно LKF составляет 17.0\% для линейного движения и 23.3\% для движения с ускорением. OC-SORT близок к EKF в линейных сценариях, но уступает при ускорении.

\subsection{Результат работы программы по отслеживанию объектов}

В качетсве проверки работоспособности программы, на вход были поданы различные видеопотоки, на которых было заснято движение транспортных средств в различных средах. Стоп-кадры с данных видеопотоков представлены на рисунках~\ref{img:deepsort_1}-\ref{img:deepsort_3}.

\img[!h]{img:deepsort_1}{deepsort_work_1.jpg}{Стоп-кадр №1 с входного видеопотока при плохом освещении}{0.75}
\img[!h]{img:deepsort_2}{deepsort_work_2.jpg}{Стоп-кадр №2 с входного видеопотока при плохом освещении}{0.75}

На изображениях видно, что у каждого зафиксированного объекта есть собственный идентификационных номер, который присваивается в соответствии с логикой работы алгоритма. Данный идентификационных номер, как можно заметить, отслеживает объект до "смерти" его трека (полное отслеживание объекта \lstinline:'21':) и при его перекрытии другими объектами (номер  \lstinline:'31':).

Как было описано в предыдущих работах, погодные условия напрямую влияют на качество детекции. Данный момент можно наблюдать на объекте с идентификационным номером \lstinline:'27':, который указан как объект \lstinline:'track':, но в реальности представляет из себя обект \lstinline:'car':.

\img[!h]{img:deepsort_3}{deepsort_work_3.jpg}{Стоп-кадр с входного видеопотока с ложной детекцией}{0.75}

Также, в некоторых случаях, наблюдается ложное срабатываение на очень похожих объектах, как, например, объект номер \lstinline:'5':, который изображен на рисунке~\ref{img:deepsort_3}.

Данные проблемы решаются путем дополнительной тренировки нейронной сети, отвечающей за определение объектов, но на работе программы по отслеживанию данная проблема никаким образом не сказывается. 